\documentclass[11pt]{article}
% Shared print layout. The web's provisional design and MathJax are independent.
\usepackage[utf8]{inputenc}
\usepackage[T1]{fontenc}
\usepackage[spanish,es-nodecimaldot]{babel}
\usepackage{lmodern}
\usepackage{amsmath,amssymb}
\usepackage{graphicx}
\usepackage{geometry}
\geometry{a4paper,left=20mm,right=20mm,top=25mm,bottom=22mm,headheight=16pt,headsep=8mm}
\usepackage{fancyhdr,lastpage}
\usepackage{enumitem}
\usepackage{xcolor}
\usepackage{pgfplots}
\pgfplotsset{compat=1.18}
\usepgfplotslibrary{fillbetween}
\usepackage[hidelinks,pdfencoding=auto]{hyperref}
\hypersetup{pdflang=es-ES,pdfauthor={},pdfsubject={Matemáticas Aplicadas a las Ciencias Sociales II}}
\pagestyle{fancy}
\fancyhf{}
\fancyhead[L]{\small Matemáticas Aplicadas a las Ciencias Sociales II}
\fancyhead[R]{\small 2.\textsuperscript{o} Bachillerato}
\fancyfoot[L]{\footnotesize Examen final | Edición revisada 1.0}
\fancyfoot[R]{\footnotesize Página \thepage\ de \pageref{LastPage}}
\renewcommand{\headrulewidth}{0.4pt}
\setlength{\parindent}{0pt}
\setlength{\parskip}{0.5em}
\setlist{itemsep=0.3em,topsep=0.3em}
\newcommand{\apartado}[1]{\subsection*{#1}}
\pgfplotsset{
  revision/.style={
    width=0.94\linewidth,height=7.5cm,
    axis lines=middle,grid=major,grid style={gray!18},
    tick label style={font=\small},label style={font=\small},
    xlabel={$x$},ylabel={$y$},samples=160,
    legend style={at={(0.5,-0.16)},anchor=north,draw=none,font=\small},
    clip=true
  }
}

\hypersetup{pdftitle={Examen final - Examen revisado}}
\begin{document}
{\Large\bfseries Examen final}

Duración: 55 minutos (normalización de la errata original). Puntuación: 10 puntos. Se permite calculadora. Se conservan los tiempos orientativos 30, 10, 5 y 10 minutos.

Apellidos y nombre: \hrulefill

Fecha: \hrulefill


\subsection*{Ejercicio 1 \hfill 5 puntos} \textbf{(30')} Se considera la función \( f(x) \) definida por:
\[
f(x) = 
\begin{cases}
ax^2 + bx + 2 & \text{si } x \leq 1, \\
\displaystyle\frac{4}{x+1} & \text{si } x > 1,
\end{cases}
\]
con \( a \) y \( b \) números reales.

\begin{enumerate}
    \item[\textbf{a)}] \textbf{(1.5 p)} \textbf{(7')} Calcule \( a \) y \( b \) para que la función \( f \) sea \textbf{continua y derivable}.    
    \item[\textbf{b)}] \textbf{(1.5 p)} \textbf{(10')} Para \( a = -1 \) y \( b = 1 \), estudie \textbf{monotonía y curvatura} de la función \( f \).
    \item[\textbf{c)}] \textbf{(0.5 p)} \textbf{(3')} Para \( a = -1 \) y \( b = 1 \), realice un \textbf{esbozo} de la gráfica de la función \( f \).
    \item[\textbf{d)}] \textbf{(1.5 p)} \textbf{(10')} Para \( a = -1 \) y \( b = 1 \), halle el \textbf{área} del recinto acotado, limitado por la gráfica de \( f \), la recta \( x = 2 \) y el eje \( OX \) (recuerda que se trata de una función definida a trozos).
\end{enumerate}

 \subsection*{Ejercicio 2 \hfill 2 puntos}\textbf{(10')} Calcule la integral de cada una de las siguientes funciones. En ambas se propone un posible cambio de variable que \textbf{no es obligatorio} realizar para resolverlas.
\[
f(x) = \frac{2x}{\sqrt{5x^2-7}},~~(t=5x^2-7); \qquad g(x) = \frac{x^2-2}{x^3-6x+5},~~ (t=x^3-6x+5).
\]

\clearpage
\subsection*{Ejercicio 3 \hfill 1 puntos} \textbf{(5')} Determine la ecuación de la \textbf{recta tangente} a la gráfica de la función \( h(x) = \displaystyle \frac{1}{x} \) en el punto de abscisa \( x = 1 \). \textbf{\textquestiondown{}La función $h$ tiene alguna asíntota? Si es así indica cuál o cuáles.}

\subsection*{Ejercicio 4 \hfill 2 puntos} \textbf{(10')} Una empresa quiere invertir en productos financieros un mínimo de un millón de euros y un máximo de seis millones de euros. La rentabilidad que obtiene viene dada en función de la cantidad invertida, \( x \), por la siguiente expresión:

\[
R(x) =
\begin{cases}
x - 2 & \text{si } 1 \leq x < 2, \\
-x^2 + 10x - 16 & \text{si } 2 \leq x \leq 6,
\end{cases}
\]

donde tanto \( x \), como \( R(x) \), están expresadas en millones de euros.

\begin{enumerate}
    \item[\textbf{a)}] \textbf{(1 p)} Estudie la \textbf{continuidad} de la función \( R \).
    \item[\textbf{b)}] \textbf{(1 p)} \textquestiondown{}Qué cantidad debe invertir para obtener la \textbf{máxima} rentabilidad y a cuánto asciende ésta? \textquestiondown{}Para qué valores de \( x \) la rentabilidad es positiva?
\end{enumerate}
















\clearpage\section*{Anexo no evaluable: acciones}
Ejercicio comentado en el original, conservado fuera de los 10 puntos activos. Su rótulo comentado indica 1 punto y sus apartados suman 2,5: esa puntuación no se adopta para evaluar este anexo.
Se desea analizar el valor de las acciones de una empresa en un día. La función \(v(t)\) nos indica el valor, en euros, de cada acción de la empresa en función del tiempo \(t\), medido en horas, a partir de la hora de apertura del mercado. De la función \(v(t)\) se conoce que su variación instantánea es:

\[
v'(t) = t^2 - 5t + 6, \quad t \in [0, 6].
\]

\begin{enumerate}
   \item[(a)] (0.75 p) Determine los intervalos de crecimiento y decrecimiento de la función \(v\).
  \item[(b)] (0.75 p) Si en el momento de la apertura del mercado se conoce que \(v(0) = 10\), halle la función \(v\).
 \item[(c)] (0.5 p) Si un inversor compró 3000 de estas acciones en el instante \(t = 2\) y posteriormente las vendió en el instante \(t = 4\), indique a cuánto ascendió la ganancia o la pérdida que obtuvo el inversor con esta gestión.
\item[(d)] (0.5 p) \textquestiondown{}En qué momentos debería haber realizado este inversor las gestiones de compra y de venta para que la ganancia hubiese sido máxima? Justifique su respuesta.
\end{enumerate}
\end{document}
